\chapter{Adversarial Robustness}
\label{ch:adversarial}

The previous chapters established that the binary random forest is the
most temporally robust flow-based detector among the four models. This
chapter asks a different question: how robust is that detector against
an attacker who knows it is there and tries to evade it? The answer is
unsatisfying. Even a simple, score-query black-box greedy attack flips
a substantial fraction of attack flows to benign within a small
perturbation budget, the resulting adversarial flows transfer almost
completely to the second-best model, and a naïve adversarial-training
defence makes the situation worse rather than better.

\section{Threat Model and Attack}

The threat model assumes a \emph{score-query black-box} attacker: the
adversary can request the model's predicted probability for any input
flow but has no access to gradients, model parameters, or training data,
and we further assume no rate limit on queries. This is the realistic
capability for an external attacker that interacts with a deployed IDS
only through the network traffic it generates. The adversary's goal is
\textit{evasion}: flipping a true-attack flow to be classified as
benign, while keeping the perturbation small enough that the resulting
flow is still operationally plausible. Of the 73 features that survive
leakage filtering, only 19 are perturbable in practice --- timing
statistics, packet lengths, packet counts, and flow duration. The
remaining 54 features are determined by the underlying network stack
and cannot be freely manipulated by an application-level attacker.

The score-query black-box assumption is realistic in the sense documented
by Apruzzese et al.~\cite{Apruzzese2023}: real-world attackers ``do not
compute gradients'' against deployed IDS, so adversarial-ML evaluations
that assume full white-box gradient access overstate the threat. A more
sophisticated robustness audit would use AutoAttack-style
ensembles~\cite{Croce2020}, but greedy coordinate-ascent is a fair upper
bound on what a black-box adversary can achieve under these constraints.

The attack itself is greedy coordinate-ascent. At each iteration, the
attacker tries six candidate step magnitudes in each of the 19
perturbable features and keeps the perturbation that most reduces the
model's predicted attack probability. The procedure runs for up to 15
iterations per flow, and the perturbation is bounded in the $L^\infty$
norm by a budget $\varepsilon$ measured in units of benign-feature
standard deviations $\sigma$.

A second important caveat: the perturbations crafted here live in
\textit{feature} space, not in \textit{packet} space. Pierazzi et
al.~\cite{Pierazzi2020} formalised this distinction as the
``problem-space'' versus ``feature-space'' gap: a flow vector that the
attack algorithm produces may not correspond to any realisable packet
sequence that the attacker can transmit on the wire. This thesis stays
in feature space; bridging the gap to genuine packet-level adversarial
flows is open work.

\section{Evasion, Transferability, and the Failure of Naïve Defence}

Two related but distinct quantities are reported. The \emph{evasion
rate} of 0.244 is measured on a sample of 500 correctly-classified
attack flows when the greedy attack is allowed to run until its
perturbation is minimised or its iteration budget is exhausted; the
median $L^\infty$ distance needed to flip a prediction is $0.250\sigma$.
The \emph{robust accuracy} curve in Table~\ref{tab:robust-acc} measures
a different quantity: the fraction of attack flows whose prediction
survives when the perturbation is hard-bounded at a fixed $\varepsilon$.
At $\varepsilon = 0.25\sigma$, robust accuracy is 0.786, meaning that
21.4\% of attack flows lose their attack label under a $0.25\sigma$
budget. The 24.4\% headline evasion rate is the looser, unbounded
version of the same effect. The curve flattens quickly after
$\varepsilon = 0.5\sigma$: most flows that can be flipped at all are
already flipped at the smallest budget.

\begin{table}[H]
    \centering
    \caption{Robust accuracy: fraction of attack flows surviving evasion at increasing budgets.}
    \label{tab:robust-acc}
    \small
    \begin{tabular}{lcccc}
        \toprule
        \textbf{Model} & \textbf{$\varepsilon = 0.25\sigma$} & \textbf{$\varepsilon = 0.5\sigma$} & \textbf{$\varepsilon = 1.0\sigma$} & \textbf{$\varepsilon = 2.0\sigma$} \\
        \midrule
        Undefended RF & 0.786 & 0.772 & 0.756 & 0.756 \\
        Adv-trained RF & 0.736 & 0.732 & 0.732 & 0.732 \\
        \bottomrule
    \end{tabular}
\end{table}

\begin{figure}[H]
    \centering
    \includegraphics[width=0.85\textwidth]{adversarial_robustness.png}
    \caption{Robust accuracy as a function of perturbation budget
    $\varepsilon$ for the undefended and adversarially-trained random
    forests. Both curves flatten quickly; adversarial training shifts
    the curve downward (worse clean accuracy) without lifting it at
    large $\varepsilon$.}
    \label{fig:adv-robust}
\end{figure}

The successful evasions are concentrated in a small set of features:
Flow IAT Min is moved in 83.6\% of successful attacks, Flow IAT Mean in
18.0\%, Flow IAT Max in 15.6\%, Forward IAT Min in 13.1\%, and Flow
Duration in 9.0\%. Packet-length features are moved in fewer than 3\%
of attacks. This concentration suggests that rate limiting, traffic
shaping, or protocol-level constraints on inter-packet timing may
reduce the effectiveness of timing-based evasion, but validating this
in practice would require packet-level experiments outside this
thesis's scope.

The adversarial flows do not stay specific to the random forest that
generated them. When the same crafted flows are re-scored by XGBoost,
the transferability rate is 86.4\% --- almost the entire RF evasion set
fools XGBoost as well. By comparison, clean attack flows are
misclassified by XGBoost at only 26\%. The practical implication is
that an attacker need not know which model the defender deploys: an
attack crafted against random forest transfers to XGBoost nearly
perfectly, so a parallel-ensemble defence in which two models must
agree before alerting offers very little additional protection.

Following Madry et al.~\cite{Madry2018}, we attempted a naïve form of
adversarial training as a defence: 2{,}000 RF-adversarial attack flows
were added to the binary training set, the random forest was retrained,
and the threshold was retuned via Youden's J. The result is a clear
negative: clean F1 dropped from 0.859 to 0.734 (a 12.5-point loss) and
the evasion rate increased from 24.4\% to 26.8\%. The naïve defence
makes both clean and adversarial performance worse simultaneously. This
is consistent with the broader literature~\cite{Tramer2020}, which has
documented that off-the-shelf adversarial training without curriculum
scheduling, careful retuning, and data balancing often hurts both
metrics. Reporting the failure is informative: it demonstrates that
adversarial robustness for flow-based IDS is not solved by adding noise
to the training set.

The practical implication for a deployed flow-based IDS is concrete.
The 21.4\% bounded evasion rate and the 86.4\% cross-model
transferability together mean that two complementary controls are
needed: model-level monitoring (tracking evasion-shaped flows that sit
close to the decision boundary across multiple retraining cycles) and
network-level enforcement (rate limiting, traffic shaping, or
protocol-level inter-packet-timing constraints that shrink the
19-feature attacker-controllable subset). Neither is in scope here;
both follow naturally from the empirical findings.
